%%%%%%%%%%%%%%%%%%%%%%%%%%%%%%%%%%%%%%%%%%%%%%%%%%%%%%%%%%
%   Vorlage von:
%
%   Prof. Dr. Bernhard Drabant
%   Prof. Dr. Dennis Pfisterer
%   Prof. Dr. Julian Reichwald
%
%%%%%%%%%%%%%%%%%%%%%%%%%%%%%%%%%%%%%%%%%%%%%%%%%%%%%%%%%%

%%%%%%%%%%%%%%%%%%%%%%%%%%%%%%%%%%%%%%%%%%%%%%%%%%%%%%%%%%
%	ANLEITUNG: 
%
%   1. Ersetzen Sie firmenlogo.jpg im Verzeichnis img
%   2. Passen Sie alle Stellen im Dokument an, die mit 
%      @stud markiert sind
%
%%%%%%%%%%%%%%%%%%%%%%%%%%%%%%%%%%%%%%%%%%%%%%%%%%%%%%%%%%

%%%%%%%%%%%%%%%%%%%%%%%%%%%%%%%%%%%%%%%%%%%%%%%%%%%%%%%%%%
%	ACHTUNG: 
%
%   Für das Erstellen des Literaturverzeichnisses wird das 
%   modernere Paket biblatex in Kombination mit biber 
%   verwendet -- nicht mehr das ältere BibTex!
%   Bitte stellen Sie ggf. Ihre TeX-Umgebung entsprechend 
%   ein (z.B. TeXStudio: Einstellungen --> Erzeugen --> 
%   Standard Bibliographieprogramm: biber)
%
%%%%%%%%%%%%%%%%%%%%%%%%%%%%%%%%%%%%%%%%%%%%%%%%%%%%%%%%%%

\documentclass[fontsize=12pt,BCOR=5mm,DIV=12,
               parskip=half,
               listof=entryprefix,paper=a4,toc=bibliography,toc=listof,pointlessnumbers
               %,headinclude=on,footinclude=off,
               %,plainheadtopline=false,plainheadsepline=false,plainfootsepline=false,plainfootbotline=false
               ]{scrreprt}
               
\makeindex

% (!!) Elementare Pakete, Konfigurationen und Definitionen werden geladen
% !TEX root =  master.tex
%      HYPERREF

%%%%%%%%%%%%%%%%%%%%%%%%%%%%%%%%%%%%%%%%%%%%%%%%%%%%%%%%%%
%	ANLEITUNG: 
%
% Passen Sie alle Stellen im Dokument an, die mit 
% @stud markiert sind
%
%%%%%%%%%%%%%%%%%%%%%%%%%%%%%%%%%%%%%%%%%%%%%%%%%%%%%%%%%%

\usepackage{makeidx}         % allows index generation
\usepackage{listings}	%Format Listings properly
\usepackage{lipsum}    %Blindtext
\usepackage{graphicx} % use various graphics formats
\usepackage[german]{varioref} 	% nicer references \vref
\usepackage{caption}	%better Captions
\usepackage{booktabs} %nicer Tabs
\usepackage{array}
\usepackage{chngcntr}
\usepackage[hidelinks=true]{hyperref} % keine roten Markierungen bei Links
\usepackage{fnpct} % Correct superscripts 
\usepackage[T1]{fontenc}
\usepackage[utf8]{inputenc}
\usepackage{calc} % Used for extra space below footsepline
\usepackage{acronym}
\usepackage{algorithm}
\usepackage{algpseudocode}
\usepackage{setspace}

%
% @stud
%
%	FONT SELECTION: Entweder 1) Latin Modern oder 2) Times / Helvetica
\usepackage{lmodern}             % 1) Latin modern font
%\usepackage{mathptmx}           % 2) Helvetica / Times New Roman fonts (2 lines)
%\usepackage[scaled=.92]{helvet} % 2) Helvetica / Times New Roman fonts (2 lines)

%
% @stud
%
%	LANGUAGE SETTINGS
\usepackage[ngerman]{babel} 	        % german language
\usepackage[german=quotes]{csquotes} 	% correct quoting using \enquote{}
%\usepackage[english]{babel}          % english language
%\usepackage{csquotes} 	              % correct quoting using \enquote{}

%
% @stud
%
% Uncomment the following lines to support hard URL breaks in bibliography 
%\apptocmd{\UrlBreaks}{\do\f\do\m}{}{}
%\setcounter{biburllcpenalty}{9000}% Kleinbuchstaben
%\setcounter{biburlucpenalty}{9000}% Großbuchstaben

%
% @stud
%
%	FOOTNOTES: Count footnotes over chapters
%1 \counterwithout{footnote}{chapter}

%	ACRONYMS
\makeatletter
\@ifpackagelater{acronym}{2015/03/20}
{\renewcommand*{\aclabelfont}[1]{\textbf{{\acsfont{#1}}}}}{}
\makeatother

%	LISTINGS
\renewcommand{\lstlistingname}{Quelltext} 
\renewcommand{\lstlistlistingname}{Quelltextverzeichnis}
\lstset{numbers=left,
	numberstyle=\tiny,
	captionpos=b,
	basicstyle=\ttfamily\small}

%	ALGORITHMS
\renewcommand{\listalgorithmname}{Algorithmenverzeichnis }
\floatname{algorithm}{Algorithmus}

%		PAGE HEADER / FOOTER
%	    Warning: There are some redefinitions throughout the master.tex-file!  DON'T CHANGE THESE REDEFINITIONS!
\RequirePackage[automark]{scrlayer-scrpage}
%alternatively with separation lines: \RequirePackage[automark,headsepline,footsepline]{scrlayer-scrpage}

%\renewcommand*{\pnumfont}{\upshape\sffamily}
%\renewcommand*{\headfont}{\upshape\sffamily}
%\renewcommand*{\footfont}{\upshape\sffamily}

\renewcommand{\chaptermarkformat}{}
\RedeclareSectionCommand[beforeskip=0pt]{chapter}
\clearscrheadfoot

%\ifoot[\rule{0pt}{\ht\strutbox+\dp\strutbox}DHBW Mannheim]{\rule{0pt}{\ht\strutbox+\dp\strutbox}DHBW Mannheim}
\ofoot[\rule{0pt}{\ht\strutbox+\dp\strutbox}\pagemark]{\rule{0pt}{\ht\strutbox+\dp\strutbox}\pagemark}
\ohead{\headmark}

\newcommand{\TitelDerArbeit}[1]{\def\DerTitelDerArbeit{#1}\hypersetup{pdftitle={#1}}}
\newcommand{\AutorDerArbeit}[1]{\def\DerAutorDerArbeit{#1}\hypersetup{pdfauthor={#1}}}
\newcommand{\Firma}[1]{\def\DerNameDerFirma{#1}}
\newcommand{\Kurs}[1]{\def\DieKursbezeichnung{#1}}
\newcommand{\Abteilung}[1]{\def\DerNameDerAbteilung{#1}}
\newcommand{\Studiengangsleiter}[1]{\def\DerStudiengangsleiter{#1}}
\newcommand{\WissBetreuer}[1]{\def\DerWissBetreuer{#1}}
\newcommand{\FirmenBetreuer}[1]{\def\DerFirmenBetreuer{#1}}
\newcommand{\Bearbeitungszeitraum}[1]{\def\DerBearbeitungszeitraum{#1}}
\newcommand{\Abgabedatum}[1]{\def\DasAbgabedatum{#1}}
\newcommand{\Matrikelnummer}[1]{\def\DieMatrikelnummer{#1}}
\newcommand{\Studienrichtung}[1]{\def\DieStudienrichtung{#1}}
\newcommand{\ArtDerArbeit}[1]{\def\DieArtDerArbeit{#1}}
\newcommand{\Literaturverzeichnis}{Literaturverzeichnis}

\newcommand{\settingBibFootnoteCite}{
	\setlength{\bibparsep}{\parskip}		  % Add some space between biblatex entries in the bibliography
	\addbibresource{bibliography.bib}	    % Add file bibliography.bib as biblatex resource
	\DefineBibliographyStrings{ngerman}{andothers = {{et\,al\adddot}},}
	\AdaptNoteOpt\footcite\multfootcite   % Will add  separators if footcite is called multiple consecutive times 
	\AdaptNoteOpt\autocite\multautocite   % Will add  separators if autocite is called multiple consecutive times
}

\newcommand{\setTitlepage}{
	% !TEX root =  master.tex
\begin{titlepage}
\begin{minipage}{\textwidth}
		\vspace{-2cm}
		\noindent \includegraphics[scale=0.25]{\imagedir/firmenlogo.jpg} \hfill \includegraphics{\imagedir/logo.jpg}
\end{minipage}
\vspace{1em}
%\sffamily
\begin{center}
	{\textsf{\large Duale Hochschule Baden-W\"urttemberg Mannheim}}\\[4em]
	{\textsf{\textbf{\large{\DieArtDerArbeit}arbeit}}}\\[6mm]
	{\textsf{\textbf{\Large{}\DerTitelDerArbeit}}} \\[1.5cm]
	{\textsf{\textbf{\large{}Studiengang Wirtschaftsinformatik}}\\[6mm]
	\textsf{\textbf{Studienrichtung \DieStudienrichtung}}}\vspace{10em}
	
	\begin{minipage}{\textwidth}
		\begin{tabbing}
		Wissenschaftliche(r) Betreuer(in): \hspace{0.85cm}\=\kill
		Verfasser(in): \> \DerAutorDerArbeit \\[1.5mm]
		Matrikelnummer: \> \DieMatrikelnummer \\[1.5mm]
		Firma: \> \DerNameDerFirma  \\[1.5mm]
		Abteilung: \> \DerNameDerAbteilung \\[1.5mm]
		Kurs: \> \DieKursbezeichnung \\[1.5mm]
		Studiengangsleiter: \> \DerStudiengangsleiter \\[1.5mm]
		Wissenschaftliche(r) Betreuer(in): \> \DerWissBetreuer \\[1.5mm]
		Firmenbetreuer(in): \> \DerFirmenBetreuer \\[1.5mm]
		Bearbeitungszeitraum: \> \DerBearbeitungszeitraum\\[1.5mm]
%		alternativ:\\[1.5mm]
%		Eingereicht: \> \DasAbgabedatum	
		\end{tabbing}
	\end{minipage}
\end{center}
\end{titlepage}
	\pagenumbering{roman} % Römische Seitennummerierung
	\normalfont	
}

%
% @stud
%
\newcommand{\settingLists}{
	%	Inhaltsverzeichnis
	\tableofcontents
	%	Abbildungsverzeichnis
	%\listoffigures
	%	Tabellenverzeichnis
	%\listoftables
	%	Listingsverzeichnis / Quelltextverzeichnis
	\lstlistoflistings
	% Algorithmenverzeichnis
	%\listofalgorithms
}

\newcommand{\initializeText}{
	\clearpage
	\ihead{\chaptername~\thechapter} % Neue Header-Definition
	\pagenumbering{arabic}           % Arabische Seitenzahlen
}

\newcommand{\initializeBibliography}{
	\ihead{}
	\printbibliography[title=\Literaturverzeichnis] 
	\cleardoublepage
}

\newcommand{\initializeAppendix}{
	\appendix
	\ihead{\appendixname~\thechapter}
}



%%%%%%%%%%%%%%%%%%%%%%%%%%%%
%
% @stud
%
%	SCHRIFTART (Schrift mit oder ohne Serifen im gesamten Text) 
%
% mit Serifen
%\addtokomafont{disposition}{\rmfamily}
%\renewcommand*{\familydefault}{\rmdefault}
%
% ohne Serifen (default)
%\addtokomafont{disposition}{\sffamily}
%
%%%%%%%%%%%%%%%%%%%%%%%%%%%%

%%%%%%%%%%%%%%%%%%%%%%%%%%%%
%
% @stud
%
% PERSÖNLICHE ANGABEN (BITTE VOLLSTÄNDIG EINGEBEN zwischen den Klammern: {...})
%
\ArtDerArbeit{Projektarbeit} % "Bachelor" oder "Projekt" wählen
\TitelDerArbeit{Vorbereitungen für automatisierte Softwaretests für Berechnungen von KPIs.}
\AutorDerArbeit{Silvan Biewald}
\Abteilung{Technology Solutions}
\Firma{Accenture}
\Kurs{WWI18DSB}
\Studienrichtung{Data Science}
\Matrikelnummer{1942347}
\Studiengangsleiter{Prof. Dr. Bernhard Drabant}
\WissBetreuer{Prof. Dr. Bernhard Drabant}
\FirmenBetreuer{Janina Rustaee}
\Bearbeitungszeitraum{24.07.2020 -- 16.11.2020}
\Abgabedatum{16.11.2020}
%
%%%%%%%%%%%%%%%%%%%%%%%%%%%%

%%%%%%%%%%%%%%%%%%%%%%%%%%%%
%
% @stud
%
%	BIBLIOGRAPHY (@stud: Bibliographie-Stil wählen - Position und Indizierung)
%
% Auswahl zwischen: IEEE Style, ALPHABETIC Style, HARVARD Style, AUTHOR-YEAR Style 
%
% (oder eigenen zulässigen Stil wählen) 
%

% Position des Zitats
%
\newcommand{\position}{inline} 
%\newcommand{\position}{footnote}

% Indizierung des Zitats
%
% 1) NUMERIC Style - e. g. [12]
%\newcommand{\indextype}{numeric} 
%
% 2) ALPHABETIC Style - e. g. [AB12]
%\newcommand{\indextype}{alphabetic} 
%
% 3) IEEE Style - numeric kind of style 
\newcommand{\indextype}{ieee} 
%
% 4) HARVARD Style 
%\newcommand{\indextype}{apa} 
%
% 5) CHICAGO Style 
%\newcommand{\indextype}{authoryear}
%
%%%%%%%%%%%%%%%%%%%%%%%%%%%%

\renewcommand*{\familydefault}{\sfdefault}

\usepackage[backend=biber, autocite=\position, style=\indextype]{biblatex} 	

\settingBibFootnoteCite

\newcommand{\abs}{\par\vskip 0.2cm\goodbreak\noindent}
\newcommand{\nl}{\par\noindent}
\newcommand{\mcl}[1]{\mathcal{#1}}
\newcommand{\nowrite}[1]{}
\newcommand{\NN}{{\mathbb N}}

\newcommand{\imagedir}{img}

\makeindex

\begin{document}

\setTitlepage

%%%%%%%%%%%%%%%%%%%%%%%%%%%%%%%%%%%%%%%%%%%%%%%%%%%%%%%%%%%%%%%%%%%%%%%%%%%%%%%%%%%%%%%%%%
% KAPITEL UND ANHÄNGE
%
% @stud:
%   - nicht benötigte: auskommentieren/löschen
%   - neue: bei Bedarf hinzufügen mittels input-Kommando an entsprechender Stelle einfügen
%%%%%%%%%%%%%%%%%%%%%%%%%%%%%%%%%%%%%%%%%%%%%%%%%%%%%%%%%%%%%%%%%%%%%%%%%%%%%%%%%%%%%%%%%%

%%%%%%%%%%%%%%%%%%%%%%%%%%%%%%%%%%%
% EHRENWÖRTLICHE ERKLÄRUNG
%
% @stud: ewerkl.tex bearbeiten
%
% !TEX root =  master.tex
\clearpage
\chapter*{Ehrenwörtliche Erklärung}

% Wird die folgende Zeile auskommentiert, erscheint die ehrenwörtliche
% Erklärung im Inhaltsverzeichnis.

% \addcontentsline{toc}{chapter}{Ehrenwörtliche Erklärung}
Ich versichere hiermit, dass ich die vorliegende Arbeit
 mit dem Thema: \textit{\DerTitelDerArbeit} selbstständig verfasst und keine anderen als die angegebenen Quellen und
Hilfsmittel benutzt habe. Ich versichere zudem,
dass die eingereichte elektronische Fassung mit der gedruckten Fassung übereinstimmt.

\vspace{3cm}
Ort, Datum \hfill \DerAutorDerArbeit
 
%%%%%%%%%%%%%%%%%%%%%%%%%%%%%%%%%%%

%%%%%%%%%%%%%%%%%%%%%%%%%%%%%%%%%%%
% SPERRVERMERK
%
% @stud: nondisclosurenotice.tex bearbeiten
%
%\input{nondisclosurenotice} 
%%%%%%%%%%%%%%%%%%%%%%%%%%%%%%%%%%%

%%%%%%%%%%%%%%%%%%%%%%%%%%%%%%%%%%%
%	KURZFASSUNG
%
% @stud: acknowledge.tex bearbeiten
%
%\input{acknowledge} 
%%%%%%%%%%%%%%%%%%%%%%%%%%%%%%%%%%%

%%%%%%%%%%%%%%%%%%%%%%%%%%%%%%%%%%%
%	KURZFASSUNG
%
% @stud: abstract.tex bearbeiten
%
% !TEX root =  master.tex
\chapter*{Abstract}

Ziel dieser Arbeit ist es für einen konkreten Anwendungsfall zwei verschiedene Automatisierungsmöglichkeiten eines Datenbankentests zu vergleichen, um einem Projektteam eine dieser Möglichkeiten zu empfehlen. Dabei handelt es sich bei der einen Variante um ein komplexes, bereits vorhandenes Test-Framework, das nicht genau auf diesen Anwendungsfall zugeschnitten ist, und bei der anderen um ein simples, vom Projektteam für diesen konkreten Anwendungsfall speziell entwickeltes. Da keine Möglichkeit eines Testlaufs der beiden bestanden, werden Vor- und Nachteile der zwei Varianten aufgezeigt und in Hinblick auf die Bedürfnisse des Anwendungsfalls gegenübergestellt. 
\newline

Die komplexere Variante übersteigt die simple nur im Punkt Zuverlässigkeit. Der Vergleich zeigt, dass für diesen Anwendungsfall die simplere Version dem Projekt durch das Zugeschnittensein auf denselben große Vorteile bietet und eine bessere Übersichtlichkeit über die Ergebnisse des Tests gewährt. Außerdem legt er nahe, dass die Zeitersparnis bei dieser Variante größer ist. Aus diesem Grund wird dem Projektteam empfohlen die einfachere Variante zu wählen.

 
%%%%%%%%%%%%%%%%%%%%%%%%%%%%%%%%%%%

%%%%%%%%%%%%%%%%%%%%%%%%%%%%%%%%%%%
% VERZEICHNISSE
%
% @stud: ggf. nicht benötigte Verzeichnisse auskommentieren/löschen in Def. von \settingLists in config.tex
%
\settingLists
%%%%%%%%%%%%%%%%%%%%%%%%%%%%%%%%%%%

%%%%%%%%%%%%%%%%%%%%%%%%%%%%%%%%%%%
% ABKÜRZUNGSVERZEICHNIS
%
% @stud: acronyms.tex bearbeiten
%
% !TEX root =  master.tex
\clearpage
\chapter*{Abkürzungsverzeichnis}	
\addcontentsline{toc}{chapter}{Abkürzungsverzeichnis}

\begin{acronym}[XXXXXXX]
	\acro{DBMS}{Database Management System}
	\acro{SQL}{Structured Query Language}
	\acro{KPI}{Key Performance Indicator}
	\acro{FDD}{Framework für Deequ basierte Datenbankentests}
	\acro{DBFALT}{Datenbanken Falsch Alt}
	\acro{DBFNEU}{Datenbank Falsch Neu}
	\acro{DBR}{Datenbanken Richtig}
	\acro{JSON}{Java Script Object Notation}
	\acro{SRS}{Software Requirements Specification}
	\acro{TRS}{Test Requirements Specification}
	\acro{UI}{User Interface}
	\acro{IDE}{Integrated Development Environment}
	\acro{RC}{Remote Control}
	\acro{HTML}{Hypertext Markup Language}
	\acro{API}{Application Programming Interface}
	\acro{RDD}{Resilient Distributed Data Set}
	\acro{Mllib}{Machine Learning library}
	\acro{AWS}{Amazon Web Services}
	\acro{S3}{Simple Storage Service}
	\acro{ETL}{Extract, Transform, Load}
	\acro{SES}{Simple Email Service}
	\acro{SDK}{Software Development Kit}
\end{acronym} 
%%%%%%%%%%%%%%%%%%%%%%%%%%%%%%%%%%%

\initializeText
\onehalfspacing

%%%%%%%%%%%%%%%%%%%%%%%%%%%%%%%%%%%
% KAPITEL
%
% @stud: einzelne Kapitel bearbeiten und eigene Kapitel hier einfügen
%
% Einleitung
%\input{introduction}

% mehrere Grundlagen- und Forschungs-Kapitel
%\input{chapter1}
%\input{chapter2}

% tatsächlicher Arbeit
% !TEX root =  master.tex
\chapter{Einleitung}
\nocite{*}

Das Testen ist ein wichtiger Teil eines jeden Softwareprojekts. Manuelles Testen kann dabei einen großen Zeit- und Arbeitsaufwand für ein Projekt bedeuten. Unter anderem aus diesem Grund werden Softwaretests automatisiert. Ziel dieser Arbeit ist es für den Testprozess eines Projektes zur Erstellung einer Anwendung, die \ac{KPI} berechnen und nutzerfreundlich anzeigen soll, Automatisierungsmöglichkeiten zu erkunden. Zu diesem Zweck werden für einen bestimmten Datenbankentest zwei Möglichkeiten zu dessen Automatisierung verglichen, ohne dass diese im vorliegenden Anwendungsfall eingesetzt werden müssen. Daraus soll eine Empfehlung resultieren, die das Projektteam bei der Entscheidung unterstützt, welche der beiden Varianten implementiert werden soll.
\newline


Zunächst werden detailliertere Informationen zu dem vorliegenden Projekt, sowie den in dieser Arbeit zu betrachtenden Tests gegeben. Darauf folgen grundlegende Themen zum Verständnis von Softwaretests, wie zum Beispiel der Testprozess oder einige Testarten. Im Anschluss werden Vorteile und Grenzen des automatisierten Testens beschrieben und zwei verschiedene Tools, die zur Automatisierung von Tests verwendet werden können, vorgestellt. Ein kurzer Überblick über einige während der Arbeit am vorliegenden Projekt verwendete Software schließen die theoretischen Grundlagen ab.
Anschließend werden detaillierte Informationen zum vorliegenden Projekt und dessen Testvorgängen näher beschrieben. Diese sind aufgeteilt in das Testen der Datenbank und das Testen der Anwendung. Für beide Fälle werden die bestehenden manuellen Tests, sowie die Konzepte für deren Automatisierung beleuchtet. Dabei ist eines davon ein bereits existierendes und das zweite ein vom Projektteam für diesen Anwendungsfall erstelltes System. Die Vor- und Nachteile der beiden Automatisierungsmöglichkeiten des Datenbankentests werden einzeln aufgezeigt und später miteinander verglichen. Zum Schluss erfolgt eine Empfehlung an das Projektteam, welche der beiden Varianten Anwendung finden sollte. Es wird erwartet, dass die vom Projektteam erstellte Variante in diesem Fall mehr Vorteile bietet.
% !TEX root =  master.tex
\chapter{Grundlegendes zum Projekt}
\nocite{*}

Die in dieser Arbeit zu vergleichenden Automatisierungsvarianten von Softwaretests stehen im Zusammenhang mit einem Projekt des IT-Beratungsunternehmens Accenture, an dem der Autor im Rahmen der Praxisphase seines dualen Studiums teilnahm. 
Gegenstand des Projekts ist es für einen international agierenden Kunden wichtige \ac{KPI}'s für dessen Unternehmen durch eine Applikation darzustellen. Bei diesen Kennzahlen handelt es sich um wichtige Leistungsindikatoren des Unternehmens, wie zum Beispiel die monatlich verkauften Einheiten eines Produkts, oder Veränderung des Umsatzes zwischen zwei Perioden.
\newline
Zum Start des Projektes bestanden bereits mehrere dieser Berichterstattungs-Anwendungen auf Seite des Kunden. Ziel des Projektes ist es nun ein System zu erstellen, das alle bestehenden Anwendungen zusammenfasst. Für jede Anwendung existiert dabei eine eigene Datenbasis. Unter Datenbasis wird im Folgenden die Menge an Daten verstanden, die benötigt wird, um alle in den Anforderungen definierten \ac{KPI}'s einer dieser Anwendungen korrekt berechnen zu können. Eine Datenbasis setzt sich wiederum aus ein oder mehreren Datenbanken zusammen. Die gesamte Datenbasis aller bestehenden Anwendungen kann Fehler enthalten und wird im Folgenden auch als \ac{DBFALT} bezeichnet. Es müssen die einzelnen Datenbasen zu einer Datenbank zusammengefasst und eine vollkommen neue Anwendung geschrieben werden.  Die Anwendung wurde im Zeitraum, in dem diese Arbeit erstellt wurde, fast vollständig fertiggestellt. Die Datenbank dieser Anwendung wird auch als \ac{DBFNEU} bezeichnet.
\newline


Die größten Schwierigkeiten ergaben sich dabei durch die Vielzahl an Quellen für die Datenbasis. 
Diese Quellen waren meist nicht homogen, sondern basierten auf unterschiedlichen Datenmodellen, so bestanden zum Beispiel unterschiedliche Bezeichnungen für ein und dieselbe Information. Deshalb wurde zunächst ein neues Datenmodell erstellt, dass im laufenden Prozess sukzessive implementiert wird. 
Dazu kam die Fehlerhaftigkeit der \ac{DBFALT}. So lagen unter anderem fehlende, oder falsche Einträge in den Tabellen der Datenbanken vor.  Das war ein Problem für die zu erstellende Applikation, da fehlerhafte Werte in den Quellen zu falschen \ac{KPI}'s in der Berichterstattungs-Anwendung führen. 
Zudem wurden im Zuge der Implementierung des neuen Datenmodells fast täglich Änderungen an der Datenbasis vorgenommen. Aus diesen Gründen machten sich Testvorgänge notwendig.

% !TEX root =  master.tex
\chapter{Grundlagen zu Softwaretests und zur verwendeten Software}
\nocite{*}
Dieses Kapitel enthält alle theoretischen Grundlagen, die für das Verständnis dieser Arbeit notwendig sind.
\section{Einführung in Softwaretests}
Zunächst werden die Grundlagen von Softwaretests in Bezug auf Definition, wichtige Terminologien, Gründe für das Testen, den Testprozess und einigen Testarten beschrieben.
\subsection{Testen - eine Begriffsdefinition}
Als eine allgemeine Definition für Testen bietet Hambling et al. folgende an: \begin{quote}Testen ist die systematische und methodische Untersuchung eines Arbeitsprodukts unter Verwendung einer Vielzahl von Techniken mit der ausdr{\"u}cklichen Absicht zu zeigen, dass es seinen gew{\"u}nschten oder beabsichtigten Zweck nicht erf{\"u}llt. Dies geschieht in einer Umgebung, die der Umgebung, die im Echtbetrieb eingesetzt werden soll, m{\"o}glichst nah kommt.~\cite[section WHAT TESTING IS AND WHAT TESTING DOES paragraph 2]{Hambling.2019}\end{quote}
Mithilfe von Tests wird also überprüft, ob ein Produkt allen Anforderungen entspricht, die an dieses gestellt werden, indem mögliche Fehler gefunden und aufgezeigt werden. Bei Softwaretest unterscheiden Hambling et al. unter drei verschiedenen Arten von  Fehlern. Zum einen dem Error, der einen menschlichen Fehler bei der Erstellung des Codes meint, zum anderen dem Defect, der sich aus einem Error ergibt, wenn dieser zum Beispiel in einem Code niedergeschrieben wird, und dem Failure, der aus einem Defect entsteht und das sichtbare falsche Verhalten eines Systems nach Ausführung des Codes meint. Dabei muss ein Defect nicht notwendigerweise immer einen Failure hervorrufen~\cite{Hambling.2019}. 
Testen ist außerdem immer ein geplanter, nachvollziehbarer und kontrollierbarer Prozess, bei dem das Produkt unter Bedingungen zu prüfen ist, die denen in der Liveumgebung möglichst nahe kommen. 

\subsection{Klärung wichtiger Begriffe zum Testen}
Um das Thema Testen zu behandeln, müssen vorab einige wichtige Begriffe geklärt werden. 
Bevor es zum Testen kommen kann, ist ein Verständnis für die Anforderungen notwendig, die das zu testende Produkt erfüllen muss. Dieses Wissen wird als Testbasis, oft zum Beispiel in der \ac{SRS} festgehalten. Die Testbasis enthält also Beschreibungen von Reaktionen, die das Produkt in bestimmten Situationen oder Zuständen zeigen soll. Diese Beschreibungen werden als Testbedingungen bezeichnet. 
Wenn für eine Testbedingung konkrete In- und erwartete Outputs festgelegt werden, spricht man von sogenannten Testfällen. 
Der Testvorgang dagegen bezeichnet die Durchführung von Testfällen, zu dem auch die Schaffung der Voraussetzungen für jeden einzelnen Testfall gehören.~\cite{Hambling.2019}
Im Folgenden wird unter Tests eine Menge von Testfällen verstanden. 

\subsection{Gründe für das Testen}
Das Testen kann unterschiedliche Gründe haben, wobei das Testen an sich keinen Wert hat oder Vorteil bringt. 
Erst infolge der Auswertung der Testergebnisse können Maßnahmen eingeleitet werden, die zum Beispiel für ein Unternehmen einen Wert generieren. 
Das Auftretens von Defects selbst führt unter Umständen zu Failures im System, die bei der Arbeit mit oder an dem Produkt stören können oder diese unmöglich machen. Sowohl im Entwicklungsprozess des Produktes als auch bei seiner Verwendung durch Endnuzter ist es dabei besonders wichtig, Defects möglichst früh zu finden, da dadurch Kosten, die für die Behebung aufgewendet werden müssen, reduziert werden können~\cite{Black.2014}. 
\newline


Andere Gründe für das Testen sind eher qualitativer Natur. Kunden erwarten, dass ein Produkt auf eine bestimmte Art und Weise funktioniert. Wenn aber ständig Defects auftreten, dann sind die Funktionen des Produktes eingeschränkt und die Erfahrung des Kunden mit dem Produkt verschlechtert sich. Daher ist es für Unternehmen sinnvoll eine Minimierung der Defects anzustreben. 
In anderen Situationen kann die Erhöhung der Verlässlichkeit einer Software sogar zur Minimierung von Risiken beitragen. ~\cite{Black.2014} Ein Beispiel hierfür ist ein Atomkraftwerk, da Failures in einem solchen Umfeld das Potenzial haben Menschenleben in großer Zahl in Gefahr zu bringen und deshalb unbedingt vermieden werden müssen.


\subsection{Der Testprozess}
Testen ist ein Prozess. Allerdings handelt es sich um keinen universellen Prozess, der einmal definiert und in jedem Fall genau so ausgeführt werden kann. Vielmehr sollte der Prozess auf den jeweiligen Anwendungsfall angepasst werden. An dieser Stelle wird also nur ein allgemeiner Testprozess vorgestellt. Eine mögliche Beschreibung eines Testprozesses beschreibt Hambling et al. folgendermaßen:
\newline
Während der Testplanung wird zunächst umrisshaft entschieden, was getestet werden soll, auf welche Weise gearbeitet wird und wie die Arbeitsverteilung sein soll. Dazu gehört auch die Erstellung der Abschlusskriterien, durch die genau definiert wird, ab wann das Testen als abgeschlossen anzusehen ist.  Die Planung wird in Form eines Testplans schriftlich festgehalten. 
\newline


Hambling et al. beschreiben als  nächstes die Testüberwachung und -steuerung, weil sie von diesem Punkt an kontinuierlich bis zum Ende des Testprozesses durchgeführt werden muss. Die Testüberwachung meint vor allem den Vergleich vom Ist-Zustand und dem im Testplan definierten Soll-Zustand. Die Teststeuerung geht mit der Testüberwachung einher, da sie die Einleitung von Maßnahmen zur Einhaltung des Testplans beinhaltet.
\newline


In der Testanalyse wird die Testbasis daraufhin analysiert, was getestet werden soll.  In diesem Schritt muss auch entschieden werden, was getestet werden kann, um daraus  die Testbedingungen zu entwickeln. In diesem Schritt werden auch mögliche Errors in der Testbasis identifiziert, um zu verhindern, dass diese später im Code implementiert zu Defects werden. 
\newline


Im Testdesign werden aus den gefundenen Bedingungen die Testfälle erstellt. Wichtig ist hierbei, wie schon im vorherigen Schritt bei den Testbedingungen, dass die Testfälle sehr klar und genau definiert werden, damit es bei der späteren Durchführung der Tests nicht zu falschem Testen kommt.
\newline


Die Testimplementierung kann bereits mit dem Testdesign einhergehen. Darunter versteht man alle Maßnahmen, die dazu dienen, die spätere Durchführung der einzelnen Testfälle zu ermöglichen. Dieser Schritt hängt maßgeblich von dem vorliegenden Anwendungsfall ab, so muss zum Beispiel für das automatisierte Testen zunächst ein entsprechendes Skript erstellt werden.
\newline


Im Schritt der Testausführung werden die einzelnen Testfälle manuell oder automatisiert durchgeführt und die Ergebnisse aufgezeichnet, sodass gefundene Defects später bearbeitet werden können. 
\newline


Wenn alle Testfälle durchgeführt wurden und alle Abschlusskriterien erfüllt sind, kann der Testprozess abgeschlossen werden. Es wird empfohlen alle für das Testen erstellten Dokumente mit den gewonnenen Erfahrungen  für zukünftige Projekte aufzubewahren ~\cite{Hambling.2019}.

\subsection{Statisches und dynamisches Testen}
Da jeder Softwaretest entweder statischer oder dynamischer Natur ist, soll auf diese beiden Konzepte hier näher eingegangen werden.
Das Testen von Software kann nicht nur während der Ausführung eines Programms geschehen, sondern auch davor. Dabei nennt sich das Testen davor, also ohne Ausführung eines Codes, statisches Testen. 
Da diese Art zu Testen keinen fertigen Code benötigt, können solche Tests teilweise bereits zu einem frühen Zeitpunkt im Lebenszyklus eines Produkts durchgeführt werden. Dies kann sehr vorteilhaft sein, weil möglichst frühzeitig gefundene Defects in der Regel weniger aufwendig zu beheben sind und so geringere Kosten verursachen. 
Unter dynamischen Testen versteht man alle Tests, die während der Ausführung eines Codes erfolgen. Im Gegensatz zum statischen Testen setzt das dynamische Testen die Ausführbarkeit der Software voraus ~\cite{Hambling.2019}.


\subsection{Einige Testarten}
In diesem Kapitel werden einige Möglichkeiten vorgestellt, wie getestet werden kann.

\subsubsection{White Box Tests}
Beim White Box-Testen wird das zu testende Programm nicht aus der Position des Anwenders betrachtet, sondern aus der des Entwicklers. Der Fokus liegt daher auf dem Testen des Programmcodes, um die Funktionalitäten des Programms sicherzustellen. Da hierbei der Code nicht ausgeführt wird, handelt es sich um statisches Testen.
\newline


Es gibt verschiedene Prinzipien, die angewendet werden können, um die Fehlerlosigkeit von Code möglichst gut sicherzustellen~\cite{Gopalaswamy.2009}. Da im Zuge dieser Arbeit jedoch das Testen von Codes keine wesentliche Rolle spielt, werden die Prinzipien und Konventionen des Testens von Codes an dieser Stelle nicht weiter ausgeführt. Falls trotz alledem Interesse daran besteht sei auf folgende grundlegenden Konzepte verwiesen: "Desk Checking", "Code Walkthrough", "Formal Inspection", die bei Gopalaswamy~\cite{Gopalaswamy.2009} beschrieben sind.


\subsubsection{Black Box Tests}
Im Gegensatz zum White Box Testen wird beim Black Box Testen der Code selbst nicht betrachtet. Die Software wird quasi von einer außenstehenden Position, also aus der Blickrichtung des Nutzers, getestet. 
Dabei werden neben der Funktionalität der Software auch deren Reaktionen auf falsche Eingaben durch den Nutzer überprüft.
Diese Vorgehensweise setzt kein Wissen um die genaue innere Funktionsweise der Software voraus. Der Wissensstand des Testers kann daher mit dem eines potenziellen Kunden übereinstimmen.
Da die Perspektive eines Nutzers eingenommen wird, beschränken sich Überprüfungen nicht nur auf die Funktionsweise bzw. das Funktionieren der Software, sondern schließen auch deren  Reaktion auf falsche Eingaben durch den Nutzer ein~\cite{Gopalaswamy.2009}.
\newline


Diese Tests basieren sowohl auf den im Vorfeld definierten Anforderungen an das fertige Produkt, als  auch auf impliziten Anforderungen, die  von einem Endnutzer erwartet werden könnten, aber in der Testbasis nicht explizit festgestellt wurden. Um sicherzustellen, dass auch alle wichtigen impliziten Anforderungen eingehalten werden, können diese neben den im \ac{SRS} festgehaltenen expliziten Anforderungen in der \ac{TRS} definiert werden~\cite{Gopalaswamy.2009}. 

\subsubsection{Regressionstests}
In manchen Fällen ist es notwendig dieselben Tests sehr häufig oder regelmäßig durchzuführen. Solche Tests werden als Regressionstests bezeichnet. Vor allem im Softwarebereich ist diese Art von Testen oft aufzufinden. Das liegt daran, dass an Software in vielen Fällen selbst nach Veröffentlichung des eigentlichen Produkts immer noch weiter gearbeitet wird. Jede neue Version der Software, die veröffentlicht wird, muss getestet werden. Gerade bei komplexen Systemen mit vielen Abhängigkeiten kann es vorkommen, dass das Hinzufügen oder Ändern einer Funktion zum einen Defects in dieser Funktion, zum anderen Defects an weiteren Stellen im System verursacht. Diese Problematik besteht bei jeder Änderung des Codes, also zum Beispiel auch bei der Behebung von Defects. Deshalb fokussieren sich Regressionstests darauf, bereits absolvierte Testfälle erneut durchzuführen. Das kann zum Beispiel nach jeder Behebung eines Defects geschehen~\cite{Gopalaswamy.2009}.


\section{Einführung in die Automatisierung von Softwaretests}
Bei der Automatisierung von Softwaretests handelt es sich um die \begin{quote}Entwicklung von Software zum Testen von Software.~\cite[section 16.1 paragraph 2]{Gopalaswamy.2009}\end{quote} 
Dieses Kapitel behandelt die Automatisierung von Softwaretests. Es wird aufgezeigt, warum Tests automatisiert werden, welche Herangehensweisen dabei verwendet werden können und welche Grenzen es gibt.

\subsection{Gründe für die Automatisierung}
Die manuelle Ausführung von Tests hat ihre Vorteile, auf die im nachfolgenden Kapitel noch eingegangen werden, erreicht aber auch schnell Grenzen. So hängt die Geschwindigkeit der Durchführung von Testfällen hier maßgeblich von menschlichen Fähigkeiten ab und erfordert oft eine große Anzahl an Testern. So kann ein bestimmter Testfall zwar wichtig für einen Test sein, jedoch einen so großen Aufwand erfordern, dass er von den Projektmitgliedern nicht geleistet werden kann. Um dem entgegenzuwirken, können Tests mithilfe von Software automatisiert werden.
\newline


Einer der wichtigsten Gründe für Testautomatisierung ist daher die Geschwindigkeitssteigerung. Die meisten Tests, die ein Mensch durchführen kann, können automatisiert viel schneller abgeschlossen werden. Dies bietet zum einen die Möglichkeit mehr Testfälle durchzuführen, zum anderen können auch komplexe oder sehr zeitintensive Testfälle effektiver ausgeführt werden. So können Testfälle, die manuell nicht durchführbar wären, in den Testprozess integriert werden. 
Für manche Tests genügt eine einmalige Ausführung nicht, zum Beispiel bei Regressionstests. Diese Art von Tests müssen häufig wiederholt werden. Auch das kann in kleineren Projekten für zu viel Aufwand sorgen.
\newline
Doch Automatisierung senkt nicht nur die Zeit, die pro Testfall benötigt wird, sondern gibt auch Arbeitskräfte  frei, die vorher manuell Tests durchführen mussten. Diese Personen können sich nun zum Beispiel darauf konzentrieren neue Testfälle zu erstellen.
\newline


Darüber hinaus werden menschliche Tester unweigerlich selbst Fehler machen, die das Testergebnis verfälschen. Dieses Problem tritt nicht auf, wenn das Testen von Software übernommen wird, sofern diese nicht fehlerhaft ist. Zudem ist keine Grenzen gesetzt, wie lange am Stück Testfälle ausgeführt werden können, da keine Ermüdungserscheinungen auftreten wie bei menschlichen Testern~\cite{Patton.2005}.


\subsection{Grenzen der Automatisierung}
Bei Betrachtung der positiven Aspekte von Automatisierung kann man leicht zum Schluss kommen, dass es immer von Vorteil ist Tests zu automatisieren. Das ist allerdings nicht der Fall. Auch der Testautomatisierung sind Grenzen geboten~\cite{Patton.2005}. Eine davon geht bereits aus der im Kapitel 3.2 (Einführung in die Automatisierung von Softwaretests) gegebenen Definition für Softwaretestautomatisierung hervor. Die zu testende Software wird bei einem automatisierten Test mithilfe von Software überprüft. Daraus folgt, dass Tests nur automatisierbar sind, wenn sie durch Software ausgeführt werden können. Software kann in vielen Fällen die Schritte eines Tests nacheinander durchführen, doch die kreativen, intellektuellen und intuitiven Eigenschaften eines menschlichen Testers fehlen hierbei~\cite{Bucsics.2015}. 
\newline


Es gibt Testfälle, die nicht automatisierbar sind. Gerade wenn das \ac{UI} zu testen ist, wird die manuelle Variante oft vorgezogen. Beispielhaft kann man sich hier eine Anwendung vorstellen, die auf verschiedenen Arten von Geräten für einen Nutzer einsetzbar sein soll. Dafür ist es wichtig zu überprüfen, ob das Interface auf allen Geräten gut erkennbar ist. Einem Menschen fällt in so einem Fall schnell auf, ob ein Text zu klein ist, oder ob bestimmte Hintergrundfarben die Lesbarkeit beeinträchtigen. Diese Art von Betrachtung ist für ein Programm nicht möglich~\cite{Bucsics.2015}.

\subsection{Tools zur Automatisierung}
Es existiert eine große Menge an Tools, die genutzt werden kann, um Softwaretests zu automatisieren. Welches davon zu wählen ist, hängt immer von dem vorliegenden Anwendungsfall ab. An dieser Stelle sollen nur die Tools näher beschrieben werden, die in dem dieser Arbeit zugrunde liegenden Projekt auch tatsächlich Anwendung fanden.
\subsubsection{Selenium}
Selenium ist ein Open Source Tool zur Automatisierung von Web Browsern. Das bedeutet, es bietet die Möglichkeit in einem Browser Maus- und Tastaturinputs zu simulieren. Der Einsatz beschränkt sich somit nicht nur auf das Testen, jedoch ist die Testautomatisierung in Browsern in den meisten Fällen das Ziel. Selenium ist ein weit verbreitetes Tool. Das liegt unter anderem an einfacher Nutzbarkeit und der Kompatibilität mit den meisten Browsern und einer Vielzahl von Programmiersprachen. Selenium ist ein übergreifendes Projekt für mehrere Tools zur Browser-Automatisierung~\cite{Axelrod.2018}.
\newline

\begin{itemize}
	
	\item Selenium WebDriver
	
	\item Selenium \ac{IDE}
	
	\item Selenium Grid
	
	\item Selenium \ac{RC}
	
\end{itemize}

Wenn von Selenium gesprochen wird, ist in den meisten Fällen Selenium WebDriver gemeint. WebDriver, auch Selenium 2 genannt, besteht zum einen aus einer Code Bibliothek, die in verschiedenen Programmiersprachen verwendbar ist, und zum anderen aus browserspezifischen Treibern. 
\newline
Mithilfe von Selenium WebDriver kann in Programmiersprachen, wie Java, Python und C\#, Code geschrieben werden, der Tastatur- und Mausinputs im Browser simuliert. Dafür muss auf die \ac{HTML} Elemente der Website zugegriffen werden. Der entsprechende Treiber führt dann beim Ausführen des geschriebenen Programms die im Code definierten Schritte durch~\cite{selenium.webdriver}.
\newline


Ein weiteres zu Selenium gehöriges Tool ist Selenium \ac{IDE}. Dies ist als ein Firefox Add-On und eine Chrome Extension umgesetzt. Eine der wichtigsten Funktionen von Selenium \ac{IDE} ist das Aufnehmen und Wiedergeben von Interaktionen mit dem Browser. Somit können Arbeitsschritte einfach und schnell automatisiert werden. Für umfangreiche Testprozesse mit einer Vielzahl von Testfällen wird es jedoch nicht empfohlen. In solchen Fällen sollte Selenium WebDriver oder Grid verwendet werden~\cite{selenium.ide}.
\newline


Um die gleichzeitige Ausführung von Testfällen auf mehreren Maschinen mit unterschiedlichen Betriebssystemen und Browsern zu koordinieren kann Selenium Grid eingesetzt werden. Einer der größten Vorteile von Grid liegt in der Zentralisierung des Testprozesses. Der Nutzer hat die Möglichkeit in seinem System unterschiedliche Konfigurationen von zum Beispiel Browserversionen für eine Reihe von Tests festzulegen. Selenium Grid sendet dann Befehle an die entfernten Browser. So können die Tests in unterschiedlichen Umgebungen parallel durchgeführt werden~\cite{selenium.grid}.
\newline

Selenium \ac{RC}, auch bekannt als Selenium 1, war das Hauptprojekt von Selenium, bevor es durch Selenium 2 ersetzt, beziehungsweise in Selenium 2 integriert wurde. In Selenium \ac{RC} wurden erstellte Skripte über den Selenium Server im Browser ausgeführt. Heutzutage wird Selenium RC nicht mehr unterstützt~\cite{selenium.rc}.


\subsubsection{Apache Spark}
Apache Spark ist nicht wie Selenium ein Tool mit dem primären Nutzen der Testautomatisierung, sondern ein Open Source Framework für parralele Datenverarbeitung. Spark ist darauf ausgelegt große Datenmengen effizient zu verarbeiten. Die schnellen Ausführungszeiten entstehen vor allem durch die Speicherung der für das jeweilige Programm benötigten Datenmengen im Arbeitsspeicher, sodass diese für das Programm jederzeit schnell abrufbar sind. Spark wurde ursprünglich in Scala implementiert, bietet aber eine \ac{API} für die Programmiersprachen Java, R, Scala und Python an~\cite{Bengfort.2016}.
Grundsätzlich bestehen Spark Anwendungen immer aus einem Treiber und einer Reihe an Arbeitern~\cite{Chellappan.2018}.
\newline

\paragraph{Spark Komponenten}
Spark setzt sich aus mehreren Komponenten zusammen.
\newline
Spark Core ist meist bekannt unter \ac{RDD} (deutsch: belastbarer verteilter Datensatz). \ac{RDD}'s sind eine Datenstruktur in Spark, die eine verteilte Menge an Objekten darstellt. Wichtig zu beachten ist, dass diese \ac{RDD}'s immutable - also nicht veränderbar sind. Das bedeutet, sollte eine Änderung an einem \ac{RDD} vorgenommen werden, dann wird dieser nicht tatsächlich abgeändert, sondern ein neuer RDD erstellt, in dem die gewünschten Änderungen zur Zeit seiner Anfertigung angewendet werden. Außerdem ermöglichen \ac{RDD}'s parallelisierte Berechnungen auf den in ihnen enthaltenen Objekten~\cite{Chellappan.2018}.
\newline
Mithilfe von Spark \ac{SQL} ist es möglich Datensätze in einer ähnlichen Weise zu manipulieren, wie es mit \ac{SQL} getan werden kann. Dafür werden als Datenstruktur hauptsächlich so genannte Dataframes verwendet. Diese besitzen im Vergleich zu \ac{RDD}'s nähere Informationen bezüglich der Struktur der Daten~\cite{Bengfort.2016}.
\newline
Eine Erweiterung von Spark Core stellt Spark Streaming dar. Durch diese Komponente ist in Spark nicht nur die sequentielle Verarbeitung großer Datenmengen möglich, sondern auch die Echtzeitverarbeitung von Streaming-Daten.
\newline
Die Spark \ac{MLib} Bibliothek besteht aus einer Reihe von Funktionen zur Unterstützung von Machine Learning für Apache Spark~\cite{Chellappan.2018}.
\newline
GraphX eröffnet die Möglichkeit die parallele Datenverarbeitung auf Graphen auszudehnen. Es besteht aus einer Reihe von Algorithmen zur Manipulation der Graphen~\cite{Bengfort.2016}.


\paragraph{Deequ}
Eine Erweiterung der Funktionalitäten von Apache Spark stellt Deequ dar. Dies ist eine von Amazon entwickelte Bibliothek zur Durchführung von Modultests auf Datenbanken~\cite{Schelter.2018}. Mit Modultests ist das separate Testen der einzelnen Komponenten eines Systems gegen deren komponentenspezifische Anforderungen gemeint~\cite{Mili.2015}. 
\newline


Ein Nutzer kann mithilfe der Funktionen von Deequ Einschränkungen und Regeln für die zu überprüfenden Tabellen festlegen. Diese werden bei der Ausführung in \ac{SQL}-Anfragen umgewandelt und durch die Unterstützung von Spark \ac{SQL} auf die Tabellen angewendet. Informationen darüber, ob die einzelnen Tests erfolgreich waren oder nicht können hinterher ausgegeben werden~\cite{Schelter.2018}.
\newline


Eine Besonderheit von Deequ ist die Fähigkeit zu testende Einschränkungen für eine gegebene Tabelle automatisch zu generieren. Dabei handelt es sich um Empfehlungen, die von Deequ nach der Inspektion einer Tabelle an einen Nutzer in Form von Deequ Funktionen, gemacht werden können~\cite{Schelter.2018}.

\section{Grundlegende Konzepte von Datenbanken}
Weil in dieser Projektarbeit das Testen von Datenbanken im Fokus steht, sollen in diesem Kapitel kurz einige grundlegende Konzepte von Datenbanken erläutert werden.
\newline

Unter einer Datenbank wird ein strukturierter Datenbestand verstanden, der durch ein \ac{DBMS} kontrolliert wird. Ein \ac{DBMS} ist eine Software, die zur Verwaltung von Datenbanken eingesetzt werden kann. Die Kombination einer Datenbank und eines \ac{DBMS} nennt man Datenbanksystem. Über das \ac{DBMS} werden Datenbanken, Datenbankmodelle zugeordnet. Diese beschreiben die Daten. Eines der meistverwendeten Datenbankmodelle ist das Relationenmodell~\cite{Saake.2018}. 
\newline

Relationale Datenbanken definieren sich darüber, dass die Inhalte einer Datenbank in Tabellen eingeteilt sind. Diese Tabellen werden auch Relationen genannt. Die Spalten werden als Attribute, und die Zeilen als Tupel bezeichnet. In dieser Arbeit werden die Begriffe Relation und Tabelle, Spalte und Attribut, sowie Zeile, Reihe und Tupel jeweils synonym verwendet. Alle in dieser Arbeit betrachteten Datenbanken folgen einem relationalen Datenbankenmodell. Um auf Datenbanken zuzugreifen wird in den meisten Fällen eine Version der Sprache \ac{SQL} verwendet~\cite{Saake.2018}.


\section{Verwendete Software}
In dem dieser Projektarbeit zugrundeliegenden Projekt wurde Software von Drittanbietern unter anderem zur Unterstützung des Testprozesses verwendet. Dieses Kapitel soll diese Software und deren für diese Arbeit relevanten Funktionalitäten beschreiben.

\subsection{Jira und Xray}
Jira ist eine Webanwendung die Teams beim Projektmanagement unterstützen kann. Zum Beispiel ist es möglich Aufgaben zu definieren, zu kategorisieren und Personen zuzuweisen. Diese werden auch als ``issues`` bezeichnet. Ein großer Fokus liegt auf agilen Arbeitsweisen, so wird unter anderem ein Kanbanboard zur Verfügung gestellt~\cite{Atlassian.}.
\newline
Gleichzeitig bietet Jira auch eine Erleichterung im Bereich des Testmanagements. Xray ist ein Add-on für Jira und ermöglicht es Tests zu definieren und zu verfolgen. Innerhalb eines Tests können eine Beschreibung des Tests, die zu schaffenden Vorherbedingungen des Tests und die einzelnen Testschritte spezifiziert werden. Für jeden Testschritt ist es zusätzlich möglich ein erwartetes Ergebnis und weitere Informationen oder Dokumente hinzuzufügen. Außerdem ist es möglich Tests zu sogenannten Testsets und auch Testplänen hinzuzufügen, was vor allem der Strukturierung dient.
\newline
Handelt es sich um einen manuellen Test, so kann dieser vom Tester gestartet und die Ergebnisse des Tests eingetragen werden. Dabei können den einzelnen Schritten Status zugeordnet werden. Alle zurückliegenden Ausführungen eines Tests bleiben gespeichert~\cite{Xray.}.

\subsection{Kurzer Exkurs zu \ac{AWS}}
\ac{AWS} ist ein Anbieter von über 175 cloudbasierten Services. Weltweit ist es die meistgenutzte Cloud Platform~\cite{Amazon.}. Im Folgenden werden nur die während des Testprozesses im vorliegenden Projekt eingesetzten Services und deren Möglichkeiten beschrieben. Hierfür wird nur ein grober Überblick über einige bestehende Funktionen dieser Services, aber keine genaueren Erklärungen zu deren Funktionsweise gegeben. Einige dieser Angebote können kostenlos sein, doch für dieses Projekt wurden nur kostenpflichtige Versionen verwendet.

\subsubsection{\ac{S3}}
Der \ac{AWS} Service \ac{S3} ist eine der meist genutzten Funktionen von \ac{AWS} und bietet die Möglichkeiten Daten in der Cloud zu speichern. Daten können nur in so genannten Buckets gespeichert werden. Um etwas in die Cloud hochzuladen muss ein Nutzer also zuerst mindestens einen Bucket erstellen. In diesen Buckets können wiederum Ordner oder Dokumente gesichert werden~\cite{Amazon.b}.
\subsubsection{Athena}
Athena ist ein Abfragedienst für Datenbanken, der es ermöglicht leicht auf in Amazon \ac{S3} gespeicherte Datenbanken zuzugreifen. Die Abfragen sind hierfür über standard \ac{SQL} zu formulieren. Nutzern ist es somit möglich die gewünschten \ac{SQL}-Abfragen einzugeben und innerhalb von Sekunden die resultierende Tabelle zu erhalten~\cite{Amazon.c}.
\subsubsection{Glue}
Mithilfe von Amazon Glue können \ac{ETL} Prozesse ausgeführt werden. In diesem Fall ist die wichtigste Funktion der so genannte Glue Job. Ein solcher Job stellt ein Skript dar, das cloudbasiert ausgeführt werden kann. Zum Beispiel ist es hiermit möglich in \ac{S3} gespeicherte Tabellen zu transformieren. Für das schnelle Arbeiten mit großen Datenmengen wird Apache Spark eingesetzt. Glue Jobs müssen in den Programmiersprachen Python oder Scala verfasst werden.
\newline


Aufgrund der Art und Weise, wie Glue mit Daten arbeitet ist es nach einem ausgeführten Glue Job, der eine Tabelle verändert oder erstellt hat, notwendig diese zu crawlen. Crawler sind eine Funktion von Amazon Glue, die Metadaten(zum Beispiel das Schema der Daten) von Datenquellen sammelt und diese speichert. Erst wenn diese Metadaten gesammelt wurden ist es wieder möglich die entsprechende Tabelle zum Beispiel über Amazon Athena anzeigen zu lassen~\cite{Amazon.d}. 
\subsubsection{\ac{SES}}
Amazon \ac{SES} erlaubt es dem Nutzer über eine Plattform von Amazon Emails zu versenden und zu empfangen. Zu diesem Zweck ist es möglich zum Beispiel mithilfe eines entsprechenden \ac{AWS} \ac{SDK}  Skripte zu erstellen um das Versenden von Emails zu automatisieren~\cite{Amazon.e}.

\input{TestvorgängeImProjekt}
% !TEX root =  master.tex
\chapter{Vergleich der Möglichkeiten zur Testautomatisierung der Datenbankentests}
\nocite{*}
In den zurückliegenden Kapiteln wurden alle Arten beschrieben, auf die im Projekt manuell getestet, sowie alle Möglichkeiten, die zur Automatisierung dieser Tests erkundet wurden.
Im Folgenden werden die beiden Varianten zur Automatisierung der Datenbankentests verglichen. Zu diesem Zweck werden erst Vor- und Nachteile beider Tests erläutert und anschließend gegenübergestellt.

\section{Variante 1 - Verwendung eines bestehenden Frameworks}
Der größte Vorteil des Einsatzes des \ac{FDD} für die Durchführung der Datenbankentests liegt in seiner Zuverlässigkeit. Das Framework wurde in einem separaten Projekt von dedizierten Softwareentwicklern erstellt und getestet. Zudem ist es bereits in mehreren Testprozessen zum Einsatz gekommen. Aus diesen Gründen kann von einer hohen Fehlertoleranz ausgegangen werden. Aktuell sind keine nicht behobenen Fehler im System bekannt. Auf der anderen Seite bietet es jedoch eine gewisse Abhängigkeit von den Urhebern. Entscheidungen über das Hinzufügen neuer Funktionen oder Änderungen des Prozesses werden vom zum Framework gehörigen Projektteam getroffen. Es ist natürlich möglich auch von außerhalb Vorschläge zu unterbreiten, deren Umsetzung ist jedoch nicht garantiert.
\newline


Das \ac{FDD} besitzt eine Reihe verschiedener Funktionen, die den Testprozess unterstützen können. So zum Beispiel die Empfehlung von Einschränkungen, die während des Testdesigns eine Grundlage bieten kann. Die Möglichkeit Testfällen eine Bewertung ihrer Wichtigkeit zuzuordnen erleichtert bei der Auswertung der aufgetretenen Defects die Priorisierung. Informationen über fehlgeschlagene Tests sind per E-Mail erhältlich. Zudem können Testfälle sowohl über Deequ Funktionen als auch über SQL-Anfragen definiert werden, was mehr Freiraum während der Testimplementierung liefert. Dafür macht es die Komplexität des \ac{FDD} aufwendiger neue Funktionen hinzuzufügen. So ist nicht nur die Entscheidung über das Hinzufügen neuer Funktionen abhängig von einer Vielzahl an Personen, sondern auch deren Implementierung zeit- und kostenaufwendig.
\newline


Außerdem fordert das Verwenden dieses Frameworks an einigen Stellen einen Extra-Zeitaufwand. Die bestehenden SQL-Anfragen müssen in das vorgegebene Format umgewandelt, und dann in ein \ac{JSON} - oder \ac{SQL} - Dokument eingetragen werden. Bei der Erstellung neuer Testfälle müssen auch neue SQL-Anfragen kreiert werden. Die erstellten Dokumente sind in den entsprechenden \ac{S3} - Buckets hochzuladen. Erst dann kann der Glue Job ausgeführt werden. 
Außerdem müssen die erwarteten Ergebnisse der Testfälle weiterhin manuell erhoben werden, da das Framework pro Glue Job den Zugriff auf nur eine Datenbank ermöglicht. Der zeitliche Aufwand wird im Vergleich zur manuellen Variante lediglich in der Testdurchführung verringert, nicht jedoch in der Testimplementierung.
\newline


Ein weiterer Nachteil liegt in der Einarbeitungszeit, die benötigt wird, um mit diesem Framework testen zu können. Um das \ac{FDD} zu verwenden ist es mindestens nötig die Dokumentation zu lesen, um Informationen über das einzuhaltende Format der Testfälle, den Speicherort für die Testfälle und die anzupassenden Parameter in den Eigenschaften des Glue Jobs zur Verfügung zu haben. Außerdem werden \ac{SQL} - Kenntnisse vorausgesetzt.




\section{Variante 2 - Erstellung eines neuen Frameworks}
Die Stärken des vom Autor erstellten Skripts liegen hauptsächlich in seiner Simplizität. Es ist einzig und allein darauf ausgelegt zwei Tabellen aus einer oder mehreren Datenbanken zu Filtern und die Zeilenanzahl der gefilterten Tabellen zu vergleichen. Dadurch ist das Skript auch für Außenstehende leicht zu verstehen, sofern sie die benötigten Programmierkenntnisse in Python und Spark besitzen. 
Die Einfachheit des Skripts ruft aber gleichzeitig den Nachteil hervor, dass Testfälle, die nicht nach dem bestehenden Prinzip aufgebaut sind, nicht auf diese Weise automatisiert werden können. 
Zum aktuellen Zeitpunkt sind keine Testfälle geplant, die von dem bisher verwendeten Schema abweichen, jedoch ist es möglich, dass sich dies in der Zukunft ändert.
\newline


Ein anderer Vorteil besteht in der Art und Weise, wie auszuführende Testfälle implementiert werden. Es  müssen hierbei nur Paare von Spalte und zu filterndem Werten in das Skript eingetragen werden. Das bedeutet es sind keine \ac{SQL}-Kenntnisse nötig, um einen neuen Testfall einzutragen. Jedoch ist hier  die Übersichtlichkeit über alle eingetragenen Testfälle eingeschränkt.
\newline


Dass das Skript extra für den vorliegenden Anwendungsfall entwickelt wurde, bedeutet zunächst einen Mehraufwand an Arbeitszeit und -kosten. Die Kosten hierfür sind jedoch relativ gering, da in das Skript nur die nötigsten Funktionen eingebaut wurden und somit keine umfangreiche Programmierung nötig war. Außerdem wurde das Programm von einem dualen Studenten und nicht einem voll ausgebildeten Softwareentwickler konzeptioniert und entwickelt, was Kosten spart.   
Es ist möglich, dass nicht bekannte Defects bestehen, da der duale Student weniger Erfahrung mitbringt, der Code von keiner zweiten Person überprüft und auch ein strukturiertes Testen des Skriptes nicht durchgeführt wurde. 
\newline


Ein weiterer Punkt sind die Freiheiten, die das Verwenden dieses Skriptes eröffnet. Alle Entscheidungen über das Skript können vom Projektteam direkt getroffen werden. Ob eine neue Funktionalität hinzugefügt werden soll, hängt allein von den  Bedürfnissen des Projektes ab. Das macht es möglich immer nur genau die Funktionen implementiert zu haben, die direkt gebraucht werden. Das spart Aufwand und senkt die Komplexität des Programms. Zudem ist das Hinzufügen neuer Funktionen auch vom Programmieraufwand her geringer, da nur wenige Abhängigkeiten zu beachten sind. Auch das ist nur dank der Einfachheit des Skriptes möglich.
\newline


Für den Testprozess bietet das Skript an mehreren Stellen Zeiteinsparungen. Die Testfälle in den Code einzutragen benötigt pro Testfall einen nur geringen Aufwand, da die einzutragenden Informationen bereits in den bestehenden \ac{SQL}-Anfragen und Testfalldefinitionen in Xray vorhanden sind. Für zukünftige Testfälle müssen außerdem keine \ac{SQL}-Anfragen mehr erstellt werden. 
Eine weitere zeitliche Einsparung passiert in der Erhebung der erwarteten Werte. Auch dafür ist bei jedem Testfall ein zeitlicher Aufwand nötig, da die hierfür benötigten Spalte - Werte - Kombinationen zunächst in das Skript einzutragen sind. Dies hat jedoch den Vorteil, dass die erwarteten Werte danach immer automatisiert erhoben werden und gleichzeitig immer aktuell sind. 
\newline 

Die Ergebnisse des Testens sind manuell über zum Beispiel Amazon Athena abrufbar. Dafür muss nach jeder Ausführung des Glue Jobs zusätzlich ein Crawler ausgeführt werden, der Metainformationen der erstellten Ergebnistabelle bereitstellt. Dieser Crawler kann immer automatisch nach der Ausführung des Glue Jobs arbeiten. Die Ausführungszeit eines solchen Crawlers beträgt wenige Sekunden. Die Einsicht in die Ergebnisse bildet somit zusätzlichen einmaligen manuellen Aufwand um einen Crawler zu konfigurieren, jedoch gleichzeitig eine gute Übersichtlichkeit über das Resultat eines jeden Testfalls. Zudem können die Ergebnisse leicht gefiltert werden.
\newline

Ein weiterer positiver Aspekt des vom Autor erstellten Skriptes ist, dass Definition und Ausführung aller Tests und Testfälle auf zwei oder mehreren Datenbanken innerhalb eines Glue Jobs zentralisiert sind. Gleichzeitig ist es mit sehr geringem Aufwand möglich zum Beispiel die Tests der einzelnen Testfallgruppen in unterschiedliche Jobs auszulagern, um die Übersichtlichkeit zu erhöhen.
\newline
Da die Testfälle direkt im Skript definiert werden müssen, werden alle bei einer Ausführung des Jobs nicht benötigten Testfälle herausgelöscht. Wenn diese jedoch bei der nächsten Ausführung wieder zu verwenden sind, müssen sie neu erstellt werden. Um  dies zu umgehen kann auf das Löschen nicht gebrauchter Testfälle verzichtet werden, wobei dann immer alle Testfälle durchgeführt werden müssten. Das würde dazu führen, dass die Ergebnistabelle durch eine Vielzahl unnötiger Informationen unübersichtlicher wird. 
Eine Möglichkeit dieses Problem zu umgehen ist zum Beispiel die Testfälle im entsprechenden Format in Skript - unabhängigen Dokumenten zu speichern und nur in das Skript einzufügen, sollten sie auch benötigt werden.

\section{Vergleich der beiden Varianten}
Zusammenfassend ist zu sehen, dass die Variante mit dem \ac{FDD} vor allem durch Zuverlässigkeit und Menge an Funktionalitäten überzeugt. Dafür ist sie jedoch unflexibler, da die Entscheidungsgewalt über das Tool außerhalb des eigenen Projektteams liegt. 
Das vom Autor erstellt Skript bekommt seine Vorteile vor allem durch Simplizität und Kontrolle des Projektteams über das Skript. Jedoch ist die Fehleranfälligkeit bei dieser Variante höher. Ausgeglichen werden kann dies über seine einfache Natur, denn so gibt es weniger Bereiche, in denen Defects auftreten können, gleichzeitig sind sie schneller auffind- und behebbar. 
In dem komplexen und langen Skript des anderen Frameworks hingegen kann das Finden von Defects einen größeren Aufwand bedeuten. 
\newline

Die Ergebnisse der Tests werden in der 1. Variante automatisch  per E-Mail an festgelegte Personen oder Teams - Kanäle versendet. Hierbei sind auch weitere Informationen zu den fehlgeschlagenen Testfällen vorhanden. Für die zweite Variante muss nach der Ausführung des Glue Jobs manuell auf die erstellte Tabelle zugegriffen werden, um die Ergebnisse zu erlangen. Dafür wird hier eine bessere Übersichtlichkeit über alle durchgeführten Testfälle gewährt. Das ist im vorliegenden Anwendungsfall von Vorteil, da hier nicht jeder gefundene Defect einzeln behoben werden muss, sondern Strukturen in den aufgetretenen Defects herauszufinden und als Ganzes zu beheben sind. Das \ac{FDD} liefert zwar schnell alle fehlgeschlagenen Testfälle per E-Mail jedoch gibt es keinen Überblick über dieselben.  
\newline

Zwei der wichtigsten Ziele der Automatisierung sind die Zeiteinsparung und das Freisetzen von Arbeitskräften, um sie an anderen Stellen einsetzen zu können. Stellt man die beiden zu Grunde liegenden Varianten gegenüber so ist zu sehen, dass in Variante 1 die Implementierung sowohl bestehender als auch neuer Testfälle einen größeren Zeitaufwand verursacht als die gleiche Implementierung mit Variante 2. Die Ausführungszeit der beiden Glue Jobs ist nicht bekannt und wird daher als gleichwertig eingeschätzt. Die zweite Variante bietet zusätzlich Zeiteinsparungen bei der Erhebung von erwarteten Werten pro Testfall. Daraus folgt, dass Variante 2 für das Projekt zeitlich einen größeren Vorteil bietet als Variante 1.

\section{Empfehlung an das Projektteam}
Dem Projektteam wird empfohlen, das vom Autor für den konkreten Anwendungsfall erstellte Skript zur Automatisierung des Datenbankentests zu verwenden. Da das \ac{FDD} in anderen Fällen schon zum Einsatz kam, bietet es zwar mehr Funktionen und könnte zuverlässiger arbeiten, unterliegt aber dennoch der anderen Variante. Diese bietet in den vom Projektteam als wichtig befundenen Kategorien, also durch ihren einfachen Charakter, die Übersichtlichkeit der Ergebnisse und die zu erwartende größere Zeitersparnis, eindeutige Vorteile. 


% Fazit und Ausblick
% !TEX root =  master.tex
\chapter{Fazit und Ausblick}
\nocite{*}

Der im Rahmen dieser Arbeit durchgeführte Vergleich zweier Automatisierungsmöglichkeiten hat ein eindeutiges Ergebnis erbracht und führte zu einer begründeten Empfehlung an das Projektteam. Im Ergebnis wird in diesem Anwendungsfall die Verwendung der vom Autor entwickelten Variante empfohlen. Obwohl davon auszugehen ist, dass die 1. Variante durch Zuverlässigkeit punkten kann, sind für dieses Projekt die Zeiteinsparungen, die Einfachheit des Systems sowie die Übersichtlichkeit der Ergebnisse von zentraler Wichtigkeit und in allen drei Kategorien ist die zweite Variante der ersten überlegen. Gleichzeitig ist jedoch anzuraten den Mehraufwand anzunehmen, das für diesen Anwendungsfall empfohlene Skript entsprechenden Softwaretests zu unterziehen, um mögliche Fehleranfälligkeiten zu reduzieren. 
\newline


Die in dieser Arbeit verwendete Methode ist insofern kritisch zu betrachten, dass keine der beiden Varianten für den vorliegenden Anwendungsfall eingesetzt werden konnte. Daher liegen keine empirischen Daten zu Zeitersparnis und Zuverlässigkeit vor. Dennoch lassen die Ergebnisse des Vergleichs eindeutige Rückschlüsse darauf zu, welche der beiden Varianten für diesen Anwendungsfall größere Vorteile bringt.
\newline


Falls den Projektmitgliedern die Ergebnisse dieser Arbeit nicht ausreichen, um eine Entscheidung zu treffen, welches System implementiert werden soll, wäre eine empirische Untersuchung der beiden Möglichkeiten in der Live-Umgebung ein nächster Schritt.


%%%%%%%%%%%%%%%%%%%%%%%%%%%%%%%%%%%

\initializeAppendix

%%%%%%%%%%%%%%%%%%%%%%%%%%%%%%%%%%%
% ANHÄNGE
%\right) \right) 
% @stud: einzelne Anhänge bearbeiten und eigene Anhänge hier einfügen 
%
% !TEX root =  master.tex
\chapter{Anhang 1}
\lstinputlisting[language=python]{Automation.py}




%\input{appendix2}
%%%%%%%%%%%%%%%%%%%%%%%%%%%%%%%%%%%

\singlespacing

%%%%%%%%%%%%%%%%%%%%%%%%%%%%%%%%%%%
% LITERATURVERZEICHNIS
% 
% @stud: Literaturverzeichnis in Datei bibliography.bib anpassen 
%
\initializeBibliography
%%%%%%%%%%%%%%%%%%%%%%%%%%%%%%%%%%%

\addcontentsline{toc}{chapter}{Index}
\printindex

\end{document}
